\chapter{Ревизия R1 и триаж ходов H0--H2}

После отрицательного закрытия C6 предложены три маршрута: честная ревизия
без exact \(\pi^{-4}\)-теоремы, повторное открытие канала
\(1\leftrightarrow3\) и объединение reciprocal-функционала с defect-saddle.
Их статусы существенно различаются.

\section{H0: независимое воспроизведение ревизии R1}

Из замороженных train-входов
\begin{equation}
 \alpha^{-1}=137.035999177,
 \qquad
 m_e=0.51099895069\ {\rm MeV},
 \qquad
 m_\mu=105.6583755\ {\rm MeV}
\end{equation}
новая прямая реализация, не импортирующая прежние численные скрипты,
воспроизводит
\begin{align}
 S_{\rm geo}&=137.0363037758784,\\
 S_{\rm vac}
 &=S_{\rm geo}-\frac1{24S_{\rm geo}}
   -\frac1{\pi^4S_{\rm geo}^2}
 =137.0359991735224,\\
 m_\tau&=1776.8594285630\ {\rm MeV},\\
 v_{\rm S2T}&=245.9934092611\ {\rm GeV},\\
 \lambda_H&=0.1292217159851,\\
 M_H&=125.0564860390\ {\rm GeV}.
\end{align}
Остаток \(S_{\rm vac}-\alpha^{-1}\) равен
\(-3.47765\cdot10^{-9}\), а относительный остаток относительно
\(m_\tau=1776.86\ {\rm MeV}\) равен
\(-3.21599\cdot10^{-7}\).

Это подтверждает машинную воспроизводимость формул, но не меняет их
доказательный статус. \(S_{\rm geo}\) остаётся геометрическим успехом,
periodic \(1/24\) --- условной determinant-ветвью,
\(\pi^{-4}\) --- strong structural compression, \(m_\tau\) --- сильной
условной связью, а абсолютный Higgs-мост наследует эти условия.

\begin{proposition}[Внутренний переход зрелости R1]
Согласно заранее записанной шкале тома, ревизия без exact
\(\pi^{-4}\)-теоремы вместе с независимой прямой реализацией выполняет
внутренний порог
\begin{equation}
 R_{\rm sci}:4/10\longrightarrow5/10.
\end{equation}
Это повышение качества фиксации и воспроизводимости, а не появление
закрытого физического предсказания:
\begin{equation}
 N_{\rm closed\ physical}=0.
\end{equation}
Порог \(5\to6\) требует внешнего воспроизведения аудита и blind-таблицы.
\end{proposition}

\section{H1: канал \(1\leftrightarrow3\)}

Предложенный тест уже выполнен. В явном quotient-нормированном
тридцатимерном coexact-базисе оболочки \(n=3\) проекция шести leaked
Killing-образов имеет
\begin{equation}
 \Tr G_{1\to3}=80,
 \qquad
 \operatorname{rank}G_{1\to3}=6,
\end{equation}
а не ноль. Консервативный denominator-proxy равен \(80/12^2=5/9\), а не
малому остатку C6. Последующий полный same-scheme аудит не нашёл
обязательной компенсации.

\begin{nogocriterion}[Запрет повторного открытия H1]
H1 закрыт отрицательно в текущей Maxwell--ghost модели. Его допустимо
переоткрыть только при выводе нового обязательного same-spectrum сектора,
точного spectral identity либо внешне доказанной эквивалентности нового
primary quotient. Повторная проекция уже вычисленного канала не является
новым маршрутом.
\end{nogocriterion}

\section{H2: reciprocal плюс defect}

Точка \(r=1\) является минимумом
\(G(r)+G(r^{-1})\) вследствие добавленной симметризации. Ambient-аудит
доказал, что текущий Maxwell--FP спектр не несёт требуемой
\(r\leftrightarrow r^{-1}\)-дуальности. Defect-saddle также остаётся
условным на parent-derived stiffness и condensation.

Для суммы этих объектов не заданы общая размерность действия, относительная
нормировка, единый вариационный принцип и карта к
\((M_W,M_Z,\sin^2\theta_W)\). Поэтому H2 пока не является вычислимым
multi-blind gate. Назначение единичного относительного коэффициента было бы
новым скрытым весом.

Кроме того, шкала тома требует сначала пройти \(5\to6\) внешним
воспроизведением, а \(6\to7\) определён как parameter-free EW/QCD matching.
H2 не может перескочить эти ворота одним стационарным пунктом.

\section{Повторная свёртка полной coexact-башни}

Повторный расчёт
\begin{equation}
 \Delta_{\rm abs}
 =\frac{\pi^2}{2}\frac{T_{\rm coex}}{S_{\rm geo}}
  \left(1-\frac{10}{24S_{\rm geo}}\right)
 =5.4667669181\cdot10^{-7}
\end{equation}
совпадает с прежним absorption-аудитом. Относительное расхождение с
\(1/(\pi^4S_{\rm geo}^2)\) равно \(3.0407\cdot10^{-6}\), а абсолютное
расхождение равно
\begin{equation}
 1.6623\cdot10^{-12},
\end{equation}
не \(1.66\cdot10^{-11}\). Оно действительно меньше train-остатка
\(|S_{\rm vac}-\alpha^{-1}|=3.4777\cdot10^{-9}\), но сравнение двух
train-остатков не создаёт нового доказательного статуса.

Необходимый коэффициент \(N_{\rm need}=10.0099700224\) показывает
несовпадение с projector-rank \(10\) в данной интерпретации. Однако
нецелое эффективное значение само по себе не запрещено для общего
взвешенного следа; запрет относится именно к отождествлению этого числа с
\(\Tr P_{0,2}\). Более фундаментально, положительная Bessel-сумма
\(T_{\rm coex}\) является Casimir kernel, а не Euclidean winding
log-determinant. Поэтому даже точное числовое устранение рангового остатка
не закрыло бы determinant theorem без функционального моста.

\begin{remark}[Статус повторной свёртки]
Формулировка «absorption с малым ранговым остатком» является уточнённым
описанием уже существующей strong structural compression, а не повышением
её статуса. Кроме того, после ревизии R1 внутренний уровень уже равен
\(R_{\rm sci}=5/10\); путь \(4\to5\) не заблокирован C6, поскольку он пройден
через честное понижение exact-утверждения и независимое воспроизведение.
\end{remark}

\begin{protocol}[Рекомендуемая последовательность]
\begin{enumerate}
 \item заморозить и выпустить пакет R1 со статусом \(R_{\rm sci}=5/10\);
 \item подготовить claim--evidence--assumption--failure-mode таблицу для
 внешнего воспроизведения \(5\to6\);
 \item не тратить вычислительный ресурс на H1 без нового обязательного
 сектора;
 \item возвращаться к H2 только после вывода общего parent action,
 относительной нормировки и заранее объявленной blind-карты.
\end{enumerate}
\end{protocol}

Машинный протокол сохранён в
\texttt{s2t\_revision\_r1\_route\_triage\_results.json}.